\chapter{2018:Frances H. Arnold and George P. Smith and Sir Gregory P. Winter}

\label{chap:chemistry2018}

The Nobel Prize in Chemistry for the year 2018 was awarded jointly to Frances H. Arnold and George P. Smith and Sir Gregory P. Winter.

\textbf{Motivation:}

for the directed evolution of enzymes (one half, Frances H. Arnold) and for the phage display of peptides and antibodies (the other half, jointly, George P. Smith and Sir Gregory P. Winter)

\section{Frances H. Arnold}

\subsection*{Early Life and Education}

Frances H. Arnold was born in 1956 in Edgewood, Pennsylvania, USA.

\subsection*{Career and Major Research}

Arnold earned a bachelor's degree from Princeton University in 1979 and a doctorate in chemical engineering from the University of California, Berkeley, in 1985. She became a professor at the California Institute of Technology, where she pioneered directed evolution, the repeated round of random mutation and screening that steers enzymes toward desired properties. She received one half of the prize.

\subsection*{Nobel Prize-winning Work}

The work for which Frances H. Arnold was awarded the Nobel Prize in Chemistry in 2018 was described by the Nobel Committee as: "for the directed evolution of enzymes".

Directed evolution harnesses natural selection in the laboratory, generating many enzyme variants and selecting those with improved stability, activity, or entirely new functions.

\subsection*{Significance and Impact}

Her methods have produced enzymes with properties not found in nature, now used to make pharmaceuticals, fuels, and chemicals by routes that are cheaper and more environmentally friendly.

\subsection*{Later Career and Honors}

Arnold continues her research on directed evolution and protein engineering at the California Institute of Technology.

\subsection*{Legacy}

Directed evolution has become a standard industrial tool and has expanded the reactions available to synthetic chemistry.

\subsection*{Selected Publications}

"Design by Directed Evolution" (Accounts of Chemical Research, 1998)

\clearpage

\section{George P. Smith}

\subsection*{Early Life and Education}

George P. Smith was born in 1941 in Norwalk, Connecticut, USA.

\subsection*{Career and Major Research}

Smith earned a bachelor's degree from Haverford College in 1963 and a doctorate from Harvard University in 1970. He became a professor at the University of Missouri, where in 1985 he invented phage display, inserting foreign genetic material into a filamentous phage so that the encoded protein appears on the surface of the virus particle.

\subsection*{Nobel Prize-winning Work}

The work for which George P. Smith was awarded the Nobel Prize in Chemistry in 2018 was described by the Nobel Committee as: "for the phage display of peptides and antibodies".

In phage display, the protein on the phage surface is physically linked to the gene that encodes it, so that proteins with desired binding properties can be selected from huge libraries and their genes recovered. He shared the other half of the prize with Gregory Winter.

\subsection*{Significance and Impact}

Phage display made it possible to select proteins and peptides with desired properties by simple binding, a process of evolution in the laboratory.

\subsection*{Later Career and Honors}

Smith is professor emeritus at the University of Missouri, where he worked for his entire career.

\subsection*{Legacy}

Phage display became the basis of protein engineering and antibody discovery, and it is used to find molecules that bind specific targets.

\subsection*{Selected Publications}

"Filamentous Fusion Phage: Novel Expression Vectors That Display Cloned Antigens on the Virion Surface" (Science, 1985)

\clearpage

\section{Sir Gregory P. Winter}

\subsection*{Early Life and Education}

Sir Gregory P. Winter was born in 1951 in Leicester, England.

\subsection*{Career and Major Research}

Winter earned bachelor's and doctoral degrees from the University of Cambridge, receiving his doctorate in 1976 while at the MRC Laboratory of Molecular Biology. There he applied phage display to antibodies, so that antibody fragments with chosen specificities could be selected and further evolved. He shared the other half of the prize with George Smith.

\subsection*{Nobel Prize-winning Work}

The work for which Sir Gregory P. Winter was awarded the Nobel Prize in Chemistry in 2018 was described by the Nobel Committee as: "for the phage display of peptides and antibodies".

Winter's antibody phage display allowed human antibodies and antibody fragments with improved properties to be produced by selection, and it underpinned the development of the first fully human therapeutic antibody, adalimumab.

\subsection*{Significance and Impact}

His work made it possible to humanize and optimize antibodies in the laboratory, giving rise to a generation of therapeutic antibodies used in medicine.

\subsection*{Later Career and Honors}

Winter is a former deputy director of the MRC Laboratory of Molecular Biology and was a founder of the company Cambridge Antibody Technology.

\subsection*{Legacy}

Antibody phage display is a foundation of modern biopharmaceutical research and has produced many approved drugs.

\subsection*{Selected Publications}

"Making Antibodies by Phage Display Technology" (Annual Review of Immunology, 1994)

\clearpage

\section*{Chapter Summary}

The 2018 prize recognized methods that use evolution in the laboratory to create useful proteins: Arnold received one half for the directed evolution of enzymes, and Smith and Winter shared the other half for the phage display of proteins and antibodies. Together these methods are central to modern biotechnology and drug development.
